\chapter{Conclusion}
\label{ch:conclusion}
In this thesis, we propose a novel approach to the problem of autonomous detumbling of non-cooperative tumbling objects in orbit using adaptive dynamic tube model predictive control (ADTMPC).
Our approach addresses the issue of unknown parameters during the operation such as the eccentricity of the orbit and the drag constant of the target.
Our approach is based on parameter estimation using set membership identification in concurrence with optimal trajectory planning for rendezvous and docking.
We further utilize robust control-invariant tubes to prevent collision violations in spite of uncertainties.
We also suggest how this can be extended to a multi-agent solution in consideration of larger targets.
Finally, we propose the logic behind detumbling the target after docking.
We have evaluated our approach on a simulated scenario involving a single chaser performing rendezvous and docking with a tumbling object in low Earth orbit.
The results show that our approach is able to successfully rendezvous, dock, and deflect the tumbling object.

We plan to extend our work and suggest future work in the implementation of a multi-agent solution including the possibility for extended observation in different directions by different chasers permitting the collection of more information over shorter time frames.
Stricter constraints for docking as suggested in Appendix \ref{app:conical-constraints} can be introduced in order to support faster and more efficient docking.
We also suggest future work in the realm of multi-target interaction where there arises the need to interact and detumble heterogeneous targets in close proximity with each other (hence posing an imminent threat of cross-collision) and with the chasers.